\documentclass[11pt,a4paper]{article}
\usepackage[italian]{babel}
\usepackage[margin=2.5cm]{geometry}
\usepackage{parskip}
\usepackage{amsmath, amssymb, amsthm}
\usepackage{booktabs}
\usepackage{array}
\usepackage{xcolor}
\usepackage{needspace}
\usepackage{enumitem}
\usepackage{listings}
\usepackage{tikz}
\usepackage[most]{tcolorbox}
\usepackage[hidelinks]{hyperref}
\usetikzlibrary{decorations.pathreplacing, shapes.geometric, arrows.meta, positioning}

% Niente vedove/orfane: i paragrafi non lasciano righe isolate a fine/inizio pagina
\widowpenalty=10000
\clubpenalty=10000

% Elenchi più compatti
\setlist{itemsep=2pt, parsep=0pt, topsep=4pt, partopsep=0pt}

% --- Riquadri con bordo nero, senza numero (si spezzano tra due pagine) ---
% Uso: \begin{definizione} ... \end{definizione}
%      \begin{teorema}[Nome] ... \end{teorema}  (il nome è facoltativo)
\tcbset{nota/.style={enhanced jigsaw, breakable, colback=white,
  colframe=black, boxrule=0.5pt, sharp corners}}
\NewTColorBox{definizione}{}{nota,
  before upper={\textbf{Definizione.}\ }}
\NewTColorBox{teorema}{o}{nota,
  before upper={\textbf{Teorema\IfValueT{#1}{ (#1)}.}\ }}
\NewTColorBox{proposizione}{o}{nota,
  before upper={\textbf{Proposizione\IfValueT{#1}{ (#1)}.}\ }}
\NewTColorBox{proprieta}{o}{nota,
  before upper={\textbf{Proprietà\IfValueT{#1}{ (#1)}.}\ }}
\NewTColorBox{assioma}{o}{nota, boxrule=1pt,
  before upper={\textbf{Assioma\IfValueT{#1}{ (#1)}.}\ }}

% --- Ambienti semplici (senza riquadro) ---
% Il titolo sta in cima al blocco, anche se il contenuto inizia con un riquadro o
% con codice; e non resta mai da solo in fondo a una pagina.
% Uso: \begin{esempio}[Titolo facoltativo] ... \end{esempio}
\newcommand{\newrunin}[2]{%
  \NewDocumentEnvironment{#1}{o}{%
    \par\Needspace{4\baselineskip}\addvspace{\medskipamount}\noindent
    \textbf{#2}\IfValueT{##1}{ (##1)}\textbf{.}\ \ignorespaces}%
  {\par\addvspace{\medskipamount}}}
\newrunin{esempio}{Esempio}
\newrunin{esercizio}{Esercizio}
\newrunin{osservazione}{Osservazione}

% V e F nelle tavole di verità
\newcommand{\V}{\textsf{V}}
\newcommand{\F}{\textsf{F}}

% --- Codice C ---
% Uso: \begin{codice} ... \end{codice}      codice C numerato
%      \begin{comando} ... \end{comando}    comandi da terminale
%      \begin{risultato} ... \end{risultato} output di un programma
%      \lstinline|printf("ciao");|           codice dentro al testo
\lstdefinestyle{c}{
  language=C,
  basicstyle=\ttfamily\small,
  keywordstyle=\color{blue}\bfseries,
  commentstyle=\color{gray}\itshape,
  stringstyle=\color{red!70!black},
  numbers=left,
  numberstyle=\tiny\color{gray},
  numbersep=8pt,
  columns=flexible,
  keepspaces=true,
  breaklines=true,
  showstringspaces=false,
  tabsize=4,
  morekeywords={include, define},
  % lettere accentate nei commenti
  literate={à}{{\`a}}1 {è}{{\`e}}1 {é}{{\'e}}1 {ì}{{\`i}}1
           {ò}{{\`o}}1 {ù}{{\`u}}1 {È}{{\`E}}1 {À}{{\`A}}1
}
\lstset{style=c}
\newtcblisting{codice}{listing only, nota, left=6mm,
  listing options={style=c}}
\newtcblisting{comando}{listing only, nota,
  listing options={basicstyle=\ttfamily\small, numbers=none, language={},
    columns=flexible, keepspaces=true, breaklines=true}}
\newtcblisting{risultato}{listing only, enhanced jigsaw, breakable,
  colback=black!5, colframe=black, boxrule=0.5pt, sharp corners,
  listing options={basicstyle=\ttfamily\small, numbers=none, language={},
    columns=flexible, keepspaces=true, breaklines=true}}

% --- Diagrammi di flusso (simboli standard) ---
\tikzset{
  terminale/.style = {ellipse, draw, minimum width=2.5cm, minimum height=0.9cm},
  io/.style        = {trapezium, trapezium left angle=70, trapezium right angle=110,
                      draw, minimum width=2.5cm, minimum height=0.9cm},
  processo/.style  = {rectangle, draw, minimum width=2.5cm, minimum height=0.9cm},
  decisione/.style = {diamond, draw, aspect=2, minimum width=2.5cm},
  freccia/.style   = {thick, -{Stealth}}
}

\title{Calcolo combinatorio con ripetizioni}
\author{Marco Cerbella}
\date{Lezione 7 -- 12 ottobre 2026}

\begin{document}
\maketitle

\section{Disposizioni con ripetizioni}

\begin{definizione}
Una \emph{disposizione con ripetizioni} di $n$ oggetti in gruppi di $k$ è una \emph{lista ordinata} di $k$ elementi scelti tra gli $n$ oggetti.
\end{definizione}

\begin{proposizione}
Il numero di disposizioni con ripetizioni di $n$ oggetti in gruppi di $k$ è
\[
D'_{n,k} = n^k .
\]
\end{proposizione}

\begin{esempio}
Quanti numeri naturali esistono con $7$ cifre tutte dispari? Gli oggetti sono le cifre dispari $\{1, 3, 5, 7, 9\}$, cioè $n = 5$. Contiamo le disposizioni con ripetizioni di $5$ oggetti in gruppi di $7$: $D'_{5,7} = 5^7 = 78\,125$.
\end{esempio}

\section{Permutazioni con ripetizioni}

\begin{definizione}
Dati $k$ oggetti distinti $a_1, a_2, \dots, a_k$, una \emph{permutazione} di $a_1, \dots, a_k$ in cui l'oggetto $a_1$ si ripete $n_1$ volte, $a_2$ si ripete $n_2$ volte, \dots, $a_k$ si ripete $n_k$ volte è una lista di $n = n_1 + n_2 + \dots + n_k$ oggetti in cui $a_1$ compare $n_1$ volte, \dots, $a_k$ compare $n_k$ volte.
\end{definizione}

\begin{proposizione}
Il numero di permutazioni di $k$ oggetti $a_1, \dots, a_k$ in cui $a_1$ si ripete $n_1$ volte, \dots, $a_k$ si ripete $n_k$ volte è
\[
\frac{n!}{n_1!\, n_2! \cdots n_k!} .
\]
\end{proposizione}

\begin{esempio}[Anagrammi]
\leavevmode
\begin{itemize}
    \item \emph{GUFO}: sono permutazioni di $4$ simboli distinti, quindi $4! = 4 \cdot 3 \cdot 2 = 24$ anagrammi (ad esempio UFOG, OFUG).
    \item \emph{GATTO}: sono permutazioni dei $4$ simboli G, A, T, O in cui T si ripete $2$ volte (liste di $5$ lettere); gli anagrammi sono $\dfrac{5!}{2!} = 60$ (ad esempio OTTAG, GATOT).
    \item \emph{MAMMA}: $A$ si ripete $2$ volte e $M$ si ripete $3$ volte, quindi $\dfrac{5!}{2!\,3!} = \dfrac{5 \cdot 4}{2} = 10$.
\end{itemize}
\end{esempio}

\section{Combinazioni con ripetizioni}

\begin{definizione}
Una \emph{combinazione con ripetizioni} di $n$ oggetti in gruppi di $k$ è un \emph{insieme} di $k$ elementi scelti tra gli $n$ oggetti (un oggetto può essere scelto più volte).
\end{definizione}

\begin{esempio}[Frutta]
Vogliamo comprare $3$ frutti e ci sono Mele, Pere, Kiwi, Arance, Banane. Rappresentiamo una scelta con \emph{stelle} (i frutti scelti) e \emph{barre} (i separatori tra i $5$ tipi di frutto). Ad esempio, scrivendo i tipi nell'ordine Mele $|$ Pere $|$ Kiwi $|$ Arance $|$ Banane:
\begin{center}
\begin{tabular}{@{}ll@{}}
$\ast \mid \ast \mid\ \mid \ast \mid\ $ & una mela, una pera, un'arancia \\
$\ast\ast\ast \mid\ \mid\ \mid\ \mid\ $ & tre mele \\
$\ \mid\ \mid \ast\ast \mid\ \mid \ast$ & due kiwi e una banana
\end{tabular}
\end{center}
Dobbiamo allora contare le permutazioni di $2$ simboli, $\ast$ e $\mid$, con ripetizioni: $\ast$ si ripete $3$ volte e $\mid$ si ripete $4$ volte. Sono
\[
\frac{7!}{3!\,4!} = 35 .
\]
\end{esempio}

\begin{proposizione}
Il numero di combinazioni con ripetizioni di $n$ oggetti in gruppi di $k$ è dato da
\[
C'_{n,k} = \frac{(n+k-1)!}{k!\,(n-1)!} = C_{n+k-1,\,k} = \binom{n+k-1}{k} .
\]
\end{proposizione}

(Nell'esempio $n = 5$ tipi e $k = 3$ frutti: $k$ stelle e $n - 1 = 4$ barre, in tutto $n + k - 1 = 7$ simboli.)

\end{document}
